Die horizontalen Schnitte sind dadurch gekennzeichnet, dass ihre vertikale Abbildung verschwindet. In der gegebenen Situation ist die vertikale Ableitung  gemäß
[[Intervall/Riemannsche Struktur/Levi-Civita-Zusammenhang/Beispiel|Beispiel]]
gleich

\mavergleichskettedisp
{\vergleichskette
{ \nabla_\partial (f)
}
{ =} { f' + f { \frac{  g' }{  2g } }
}
{ } {
}
{ } {
}
{ } {}
}
{}{}{.}
Dafür ist zunächst

\mavergleichskette
{\vergleichskette
{ f
}
{ = }{ 0
}
{  }{
}
{    }{
}
{   }{
}
}
{}{}{}
eine Lösung. Es sei nun

\mavergleichskette
{\vergleichskette
{ f(t_0)
}
{ \neq }{ 0
}
{  }{
}
{    }{
}
{   }{
}
}
{}{}{}
und dies gilt dann auch in einer offenen Umgebung von $t_0$. Darauf ist

\mavergleichskettedisp
{\vergleichskette
{  { \frac{  f' }{ f } }
}
{ =} { { \frac{  g' }{  2g  } }
}
{ } {
}
{ } {
}
{ } {
}
}
{}{}{.}
Das Lösungsverfahren für lineare Differentialgleichungen führt auf

\mavergleichskettedisp
{\vergleichskette
{   \ln  \betrag {  f  }
}
{ =} { c+ { \frac{   1  }{  2 } }  \ln   g
}
{ } {
}
{ } {
}
{ } {
}
}
{}{}{}
und somit auf

\mavergleichskettedisp
{\vergleichskette
{ \betrag {  f  }
}
{ =} { e^c e^{ { \frac{   1  }{  2 } }  \ln  g }
}
{ =} { e^c  \sqrt{g}
}
{ } {
}
{ } {
}
}
{}{}{}
bzw.

\mavergleichskettedisp
{\vergleichskette
{ f
}
{ =} { d \sqrt{g}
}
{ } {
}
{ } {
}
{ } {
}
}
{}{}{}
\zusatzklammer {mit

\mavergleichskette
{\vergleichskette
{ d
}
{ \in }{  \R
}
{  }{
}
{    }{
}
{   }{
}
}
{}{}{}} {} {,}
was man auch direkt überprüfen kann.

 [[Kategorie:Latexseite]]
